\chapter*{Abstract}

\addcontentsline{toc}{chapter}{Abstract}
Traditional academic credential verification systems face significant limitations in addressing modern educational needs. Paper-based and centralized digital credentials are vulnerable to forgery and require direct institutional contact for verification, creating administrative burdens and delays. While blockchain-based credential systems have emerged to provide immutability and decentralized verification, existing solutions suffer from critical gaps: they can only verify whole certificates rather than individual courses, produce proof sizes that scale with dataset size, lack privacy-preserving selective disclosure, and do not provide cryptographic guarantees of issuance timing. This thesis introduces IU-MiCert, a blockchain-based system for verifiable academic credentials that addresses these limitations through advanced cryptographic tree structures. The system enables granular verification where individual courses can be validated independently, provides compact constant-size proofs that remain efficient regardless of institutional scale, supports privacy-preserving selective disclosure allowing students to reveal only relevant achievements, and establishes verifiable temporal integrity through blockchain-anchored issuance timelines. A novel revocation mechanism based on credential versioning provides institutions with the ability to invalidate credentials while maintaining system integrity. The research demonstrates that this approach achieves significant improvements in proof efficiency, verification speed, and economic viability compared to existing solutions. The complete implementation, including command-line tools, web interfaces, and deployed smart contracts, validates the practical feasibility of the approach. Economic analysis shows substantial cost reductions compared to traditional verification processes, making the system viable for institutional deployment. The work contributes to advancing blockchain credential technology by demonstrating successful application of modern cryptographic primitives to real-world credential management challenges, providing a foundation for next-generation academic verification infrastructure that balances security, privacy, efficiency, and usability.

\vspace{0.5cm}

\noindent\textbf{Keywords:} micro-credentials, academic provenance, blockchain verification, credential fraud prevention, learning timeline, verkle tree